\chapter{Pointing at the Source}
\label{chap:pointing-at-the-source}

The machine is built. Over its own values it has written, and had certified, every rung of itself, from
the first mark to the compiler tape that folds back onto what it compiles. In all of that building, the
thing being measured was always handed in from outside --- a value, a present, an event, a procedure ---
and the apparatus read it. This movement hands the apparatus the one carrier it has never been run over:
not a value from the world, but its own truth. What it will measure now is whether things are true, and the
first thing it will turn that power on is itself. This short chapter stages the move; the next springs what
the move releases.

\section{The carrier is its own truth}
\label{se:carrier-is-its-own-truth}

Recall what a carrier is: the thin process that supplies the things being told apart and the order they
fall into, off which every reading is taken. Every carrier the machine has used until now supplied values
--- counts, bits, magnitudes, events. The carrier this movement installs supplies something else. It
supplies \emph{truth-values}: its things-to-be-told-apart are propositions, and its order is the order of
truth. The distinguished symbol the apparatus reads against --- the standard it compares everything to ---
is set, for the first time, to truth itself. The apparatus \emph{tanges} by truth now: where it once
selected one value from another by a decidable difference, it selects the true from the false, reading not
how much but whether. Nothing in the tower is rebuilt to do this. The carrier is swapped, and the same
machine that read values is pointed at the question of truth.

This is a stranger move than it sounds, and it is worth dwelling on why. A measuring instrument that
measures truth-values is measuring the very currency in which its own verdicts are denominated. Every Fact
the machine has ever used was a proposition with a way to decide it; now the propositions themselves are
the carrier. The apparatus is no longer reading the world through its truth; it is reading truth as though
truth were a world. That turn --- from reading by truth to reading truth --- is the whole content of
self-hosting, and the rest of the chapter is the consequence of having made it.

\section{Re-instantiating the tower}
\label{se:re-instantiating-the-tower}

With the carrier swapped, the construction does the one thing the bootstrap was built to make possible: it
runs the entire tower again, over the new carrier, without writing anything new. Every rung --- the bare
difference, the number tower, the bit and the clock, the present and the gauge, the source and the
execution, the order on its own knowledge, the local and the universal, the halt and the measure and the
compiler tape --- is re-instantiated over truth. The admissible count becomes a count of truths; the order
becomes the order of truth-values; the compiler tape becomes a tape that compiles propositions. Thirty-five
rungs, the same thirty-five, now reading the apparatus's own truth instead of values handed in from
outside.

That this is even possible is the dividend of how the bootstrap was built. Because every rung was
constructed over an arbitrary carrier, holding for whatever values it was handed, it holds equally when the
values it is handed are truth-values. The construction does not have to argue the tower over again; the
certifier that accepted each rung over a generic carrier has, in accepting it, already accepted it over this
one. The machine that read the world is, with no new instruction, the machine that reads its own truth ---
which is precisely what it means for a compiler to be written in the language it compiles. It can take its
own source as input because its own source is written in the carrier it now runs over.

\section{The tape, handed its own source}
\label{se:tape-handed-its-own-source}

This is the self-hosting moment, and it is exactly the compiler tape of the previous chapter handed a new
input: its own source. In the language of the subject, the construction has become a self-hosting
compiler, a proof assistant pointed at the propositions it was built to check, a type-theoretic elaborator
asked to elaborate its own truth. The tape that folded back onto the values it compiled now folds back onto
the propositions that describe it; the boot, the strapping-up, and the compute run, this time, on the
machine's own account of itself. Nothing is added. The same convergence the tape was proved to have --- the
compiled object strictly refining its source --- is now convergence of the apparatus's reading of its own
truth toward that truth.

A compiler handed its own source is the sharpest test a construction of this kind can face, because there
is nowhere left to hide an unearned assumption: anything the apparatus smuggled in would have to survive
being read by the apparatus itself. That is the discipline the whole bootstrap was building toward, and it
is why the movement matters. The machine is about to read the one source it cannot fool, and what it finds
there is the subject of the next chapter.

\section{What is about to give}
\label{se:what-is-about-to-give}

It is worth naming, before the next chapter, the two concepts the reading will turn on, because they are
doing the work. The first is proof
irrelevance: at the level of truth, what matters is \emph{that} a proposition holds, not how one came to
hold it, so that any two justifications of the same truth are treated as the same --- the apparatus
\emph{funges} all the proofs of a proposition into one, pooling them by the single characteristic that they
establish the same truth. The second is propositional extensionality: two propositions that hold under
exactly the same conditions are, at this level, the very same proposition. Held together --- every
justification of a truth identified, and propositions identified when their conditions coincide --- these
are the levers the next chapter pulls.

They are levers because the apparatus has, all along, built everything on telling apart. The bare
difference was the first tange and the floor of the tower. But the carrier is now truth, and at truth, the
two principles just named begin to erase exactly the differences the machine relies on. A distinction that
cannot survive proof irrelevance and propositional extensionality is a distinction the apparatus, pointed
at its own truth, can no longer draw. The machine has aimed its sharpest faculty --- telling apart --- at
the one carrier where that faculty is about to fail. The next chapter is the failure, and why it is not a
defect but the point.

\section{What is demonstrated, and what is assumed}
\label{se:demonstrated-assumed-II1}

What is demonstrated is the move itself, and it is proved rather than merely described: the carrier is set
to the apparatus's own truth, and the entire thirty-five-rung tower is re-instantiated over it with no new
construction, because each rung was built over an arbitrary carrier and a certifier that accepted it there
has already accepted it here. The compiler tape is handed its own source --- the self-hosting moment ---
and the convergence proved for it over values carries to its reading of its own truth. Self-hosting is
staged, not yet sprung.

What is assumed is unchanged. The seed and its decidability remain the foundation; the honest import of a
real machine's behavior, at the two rungs that needed it, still stands and is still marked as borrowed.
Nothing new is granted to make the carrier-swap work --- it works because the tower was built general
enough to allow it. And the audit the previous chapter deferred is still deferred: the collapse the next
chapter proves will rest on the two principles just named, but the question of exactly what the whole
construction rests upon --- its final tally of foundations --- is not opened here and will not be opened
until the closing movement weighs it. The single standing grant, the law of motion the first book named at
its summit, remains out of play. The apparatus is pointed at its own truth, and has not yet read it.
